\subsection{UC-05: Gestão do Catálogo}

\begin{itemize}[topsep=1pt, itemsep=0pt]
  \item[]\textbf{Descrição}: Gerente da loja adiciona, visualiza, atualiza e remove pratos do catálogo.

  \item[]\textbf{Ator Principal}: Gerente.

  \item[]\textbf{Atores Secundários}: Cliente.

  \textbf{Pré-condições}:
  \begin{enumerate}
    \item Gerente autenticado;
    \item Loja gerenciada pelo Gerente existe;
    \item Sistema está disponível.
  \end{enumerate}

  \textbf{Pós-condições}:
  \begin{enumerate}
    \item Prato adicionado, atualizado ou removido do catálogo;
    \item Catálogo reflete as mudanças imediatamente;
    \item Auditoria registra todas as operações;
    \item Clientes veem o catálogo atualizado.
  \end{enumerate}

  \textbf{Fluxo Principal - Adicionar Prato}:
  \begin{enumerate}
    \item Gerente acessa a página de catálogo e preenche o formulário para adicionar um novo prato;
    \item Sistema valida os dados;
    \item Sistema retorna os detalhes do prato criado;
    \item Ação é registrada na auditoria.
  \end{enumerate}

  \textbf{Fluxo Alternativo - Listar Pratos}:
  \begin{enumerate}
    \item Gerente ou cliente acessa o catálogo;
    \item Sistema retorna lista paginada de pratos disponíveis;
    \item Para cada prato: nome, preço e disponibilidade;
  \end{enumerate}

  \textbf{Fluxo Alternativo - Visualizar Detalhes do Prato}:
  \begin{enumerate}
    \item Usuário seleciona um prato específico;
    \item Sistema retorna todas as informações do prato;
  \end{enumerate}

  \textbf{Fluxo Alternativo - Atualizar Prato}:
  \begin{enumerate}
    \item Gerente seleciona um prato para atualizar;
    \item Sistema valida os novos dados;
    \item Prato é atualizado no banco de dados;
    \item Sistema retorna os dados atualizados;
    \item Ação é registrada na auditoria com o que foi alterado.
  \end{enumerate}

  \textbf{Fluxo Alternativo - Remover Prato}:
  \begin{enumerate}
    \item Gerente seleciona um prato para remover;
    \item Sistema confirma a exclusão e processa a remoção;
    \item Sistema remove as associações de ingredientes (dish\_ingredient);
    \item Ação é registrada na auditoria.
  \end{enumerate}

  \textbf{Regras de Negócio}:
  Consulte \ref{subsec:catalogue}.
\end{itemize}
